The Pantheon+ distance moduli, DES Y6 / KiDS cosmic shear, and galaxy–galaxy lensing band-powers represent foundational, complementary datasets utilized in modern precision observational cosmology to constrain cosmic expansion and structure growth.Pantheon+ Distance ModuliThe Pantheon+ Analysis compiles over 1,700 Type Ia Supernovae (SNe Ia) light curves spanning a redshift range up to z ~ 2.3.
Core Measurement: It extracts the observed distance modulus (\(\mu_{\rm obs} = m_b - M\)), mapping the luminosity distance (\(d_{L}\)) as a function of redshift.
Systematics & Calibration: It implements the "SuperCal–Fragilistic" cross-calibration scheme, corrects for local peculiar velocity fields using 2M++ maps, and includes localized systematic uncertainties directly into its comprehensive covariance matrix.
Cosmological Role: By introducing SNe Ia within Cepheid-hosting galaxies calibrated by the SH0ES project, Pantheon+ breaks the absolute magnitude (M) degeneracy, directly constraining the local Hubble constant (H₀) alongside the dark energy equation of state (w).
DES Y6 / KiDS Weak Lensing & Band-Powers: The final Dark Energy Survey Year 6 (DES Y6) results and the Kilo-Degree Survey (KiDS) map the distribution of matter using weak gravitational lensing and large-scale structure. Instead of real-space correlation functions, these analyses often employ band-powers (Fourier-space or filtered angular power spectra) to isolate specific spatial scales and avoid complex small-scale modeling issues.
Cosmic Shear: Derived from the shapes of millions of background source galaxies (e.g., ~ 140 million in DES Y6 Metadetection), isolating matter fluctuations (\(S_8 \equiv \sigma_8 \sqrt{\Omega_{\rm m}/0.3}\)).
Galaxy–Galaxy Lensing (GGL): Measures the cross-correlation between foreground lens positions (e.g., MagLim++ sample) and background source shear to map out galaxy bias and matter distribution.
The Role of \(\mu _{\mathrm{l}ens}\) (Lensing Amplification)In joint analyses utilizing galaxy-galaxy lensing and cosmic shear, \(\mu _{\mathrm{l}ens}\) (often formalized alongside scaling parameters like \(A_{\mathrm{l}ens}\) or related magnification / shear calibrations) acts as a critical scaling parameter.
Lensing Calibration: It quantifies any systematic deviations or phenomenological modifications from the expected General Relativity lensing amplitude.
Magnification Integration: It captures weak lensing magnification effects—where foreground structures alter the observed selection thresholds, fluxes, and densities of background samples—ensuring that the underlying matter-bias relationship remains completely unbiased at high statistical precision.
